\subsection{Deeper pure radicals and ray class generators away from two and three}

\begin{statement}{A.105}{Proved}
{Compatible deeper root splitting and the coefficient-place norm map}
Retain $a=3u^8$, $b=-2v^8$, $u,v$ prime to six,
$N=3^n$, $n\ge1$, and the actual representative
$\beta\in K_n^\times$ of A.101, with its canonical
two-adic divisor and $N_{K_n/\Q}\beta=b$.
Choose $t>n$, compatible roots
\[
 \theta^{3^t}=a,\qquad
 \alpha=\theta^{3^{t-n}},\qquad K_n=\Q(\alpha),
 \qquad L_j=\Q(a^{1/3^j},\mu_{3^j})\quad(n\le j\le t).
 \tag{A.105.1}
\]
Let $F$ be the compositum of $L_n$ with the
$\Q$-conjugates of a three-primary ray class field
of $K_n$ whose finite modulus is supported away
from \emph{both} two and three.
Choose $r\ge1$, $\operatorname{lcm}(24,N)\mid M$, and
\[
 R=\Q(\mu_D),\qquad
             \operatorname{lcm}(M,3^r,3^t)\mid D.
 \tag{A.105.2}
\]
With the actual $S_\beta$ and $B_R$ of A.101, put
$S_t=S_\beta B_RL_tR$.
Then the precise intersections and abelianization are
\[
 FR\cap L_tR=L_nR,\qquad
             FR\cap S_t=L_nR,\qquad S_t^{\rm ab/\Q}=R.
 \tag{A.105.3}
\]
Thus the split-prime class generators in A.95 and
A.100, for these ray moduli, may simultaneously split
$L_t$, every displayed conjugate radical of $\beta$,
and the full boundary field $\Q(\mu_M,b^{1/M})$.
Every larger cyclotomic CRT automorphism fixing $R$
remains compatible with these conditions.
The permitted primes satisfy $v_3(q-1)\ge t$.

At the unique place of $K_n$ above three, put
$k_n=\Q_3(\alpha)$.
For every $Q\in\Q^\times$ and integer $e$ prime to three,
\[
 k'_n=k_n(\gamma),\qquad\gamma^3=Q^e\alpha,
                  \qquad N_{k'_n/k_n}(k_n'{}^\times)=k_n^\times.
 \tag{A.105.4}
\]
The norm map on their three-primary completions is
also onto.
In particular this coefficient-place norm equation
places no additional restriction on the unit of the
target representative.
\end{statement}
\begin{proof}
\emph{Total ramification of the deeper radical at three.}
Fix a completion $R_w$ at three.
It is an abelian extension of $\Q_3$.
The element $a$, of three-adic valuation one, is not
a cube in any abelian extension of $\Q_3$.
Indeed a cube root would generate an Eisenstein
cubic extension of $\Q_3$ which is not Galois:
normality would require $\mu_3$, whose degree is two.
Every subfield of an abelian extension is Galois,
giving the contradiction.
Since $R_w$ contains $\mu_{3^t}$, Kummer theory
therefore gives
\[
 [R_w(a^{1/3^t}):R_w]=3^t,
 \qquad
 \operatorname{Gal}(R_w(a^{1/3^t})/\Q_3)
       =C_{3^t}\rtimes\operatorname{Gal}(R_w/\Q_3).
 \tag{A.105.5}
\]
The splitting in (A.105.5) is literal: an automorphism
of $R_w$ fixes the selected root $a^{1/3^t}$.
The full Kummer degree makes these assignments
well-defined on the root algebra.
Its action on $C_{3^t}$ is the cyclotomic action.
An element whose cyclotomic value is $-1$ modulo three
makes the commutator subgroup contain $C_{3^t}$.
The quotient by this translation group is abelian,
so the maximal abelian subfield in (A.105.5) is exactly
$R_w$.

Every unramified quotient of a finite local Galois
extension is cyclic and hence abelian.
It follows that $R_w(a^{1/3^t})$ and $R_w$ have the
same residue field: their maximal unramified
subextensions over $\Q_3$ are the same.
Thus $L_tR/R$ is totally ramified at three,
and $L_tR/L_nR$ is totally ramified there with degree
$3^{t-n}$.
In contrast, $F/L_n$ is unramified at three by the
specified ray modulus, and remains so after base
change to $L_nR$.
An extension which is both unramified and totally
ramified is trivial, proving the first intersection
in (A.105.3).

\emph{The simultaneous radicals and the intersection at two.}
At two, the odd power $3^t$ has a unique root of $a$
in $\Q_2$.
Since $R$ contains $\mu_{3^t}$, $L_tR/R$ splits
completely at every prime above two.
At its primes over a fixed $w\mid2$ of $R$, the
matrix of valuations of the $3^n$ conjugates of
$\beta$ is the matrix of A.101 with each row
repeated $3^{t-n}$ times.
Selecting one row over each old prime gives
\[
                        e(R_w/\Q_2)I_{3^n}.
 \tag{A.105.6}
\]
This scalar is a three-adic unit.
All conjugate Kummer classes therefore remain
independent after adjoining $L_tR$, at every finite
three-power precision.

Put $m=v_3(M)$ and $h=\max(r,m)$.
After also adjoining the prime-to-three part of
$B_R$, the simultaneous three-primary Kummer lattice
is generated, modulo $3^h$th powers, by
\[
 \beta_i^{3^{h-r}} (0\le i<3^n),
                  \qquad b^{3^{h-m}}=
                            \prod_i\beta_i^{3^{h-m}}.
 \tag{A.105.7}
\]
The base change multiplies (A.105.6) only by
ramification indices prime to three.
The valuation map on this lattice is injective.
Every nonzero cyclic Kummer subextension consequently
ramifies at some prime above two, by the exact
valuation-denominator argument of A.101.
The extension $FL_tR/L_tR$ is a three-extension
unramified at all these primes.
The prime-to-three boundary part has trivial
intersection with it by degree, and the displayed
three-primary part has trivial intersection by
ramification.
It follows that
\[
                  (FL_tR)\cap S_t=L_tR.
 \tag{A.105.8}
\]
Intersecting with $FR$ and applying the first equality
of (A.105.3) proves $FR\cap S_t=L_nR$.

\emph{The actual arithmetic group and CRT compatibility.}
Let $d=\min(r,m)$.
The three-primary group over $L_tR$ is exactly
\[
 V'=\{(z_i,w)\in(\Z/3^r\Z)^{3^n}\oplus\Z/3^m\Z:
                         \sum_i z_i=w\pmod{3^d}\}.
 \tag{A.105.9}
\]
This follows from (A.105.6) and the same exact
common-$b$ relation as in A.101.
A translation lift sends
\[
 \theta\longmapsto\zeta_{3^t}\theta,
       \qquad \delta_i\longmapsto\delta_{i+1},
       \qquad b^{1/3^m}\longmapsto b^{1/3^m},
 \tag{A.105.10}
\]
where the $\delta$ indices are reduced modulo $3^n$.
It has order $3^t$.
A cyclotomic automorphism of value $u$ fixes $\theta$
and $b^{1/3^m}$ and sends $\delta_i$ to $\delta_{ui}$.
The root relations and the common-$b$ relation are
preserved, so these are actual automorphisms.
Thus the simultaneous three-primary arithmetic
kernel is $V'\rtimes C_{3^t}$ with this action.
Its translation coinvariants on $V'$ are the cyclic
group of order $3^h$ from A.101.
Cyclotomic $u=-1$ kills both these coinvariants and
the full translation $C_{3^t}$ in arithmetic
abelianization.
The prime-to-three boundary part is treated by the
actual fibre product of A.101; it has no abelian
subfield beyond $R$.
This proves $S_t^{\rm ab/\Q}=R$.

If $R'/R$ is a larger cyclotomic field, then
$S_t\cap R'=R$.
Hence the identity on $S_t$ and any CRT automorphism
of $R'$ fixing $R$ define an automorphism of $S_tR'$.
The intersection proof above remains valid after
replacing $R$ by $R'$, giving
$FR'\cap S_tR'=L_nR'$.
The Hilbert Frobenius coset $\sigma_0V_0$ of A.95
therefore remains compatible with these additional
splitting conditions.
Its quotient $V/V_0$ is still central, since it
arises from $F\cap R'$.
The identical coset-character proof gives the full
split-prime class target, and its subgroup-$J$
version gives the same-prime relative target of
A.100.
This establishes the asserted preservation of those
actual class generators.
Chebotarev supplies the primes outside any finite
exceptional set and above any fixed bound.
Complete splitting in $\mu_{3^t}$ gives
$v_3(q-1)\ge t$.

\emph{The coefficient-place norm.}
The field $k_n$ has odd degree $3^n$ over $\Q_3$
and does not contain $\mu_3$.
Its normalized valuation gives
\[
             v_{k_n}(Q^e\alpha)=1+3^n e\,v_3(Q),
 \tag{A.105.11}
\]
which is prime to three.
Thus $k'_n/k_n$ has degree three and is not Galois;
its closure $\widetilde k_n=k'_n(\mu_3)$ has group
$S_3$ and maximal abelian subfield $k_n(\mu_3)$.
The local norm limitation theorem gives
\[
 [k_n^\times:N_{\widetilde k_n/k_n}
                  (\widetilde k_n^\times)]=2;
 \tag{A.105.12}
\]
see \cite[Remark 4.3.7, equation (4.3.1), and the Galois
case of Theorem 4.1.7]{BIR_KedlayaCFT}.
The subgroup $N_{k'_n/k_n}(k_n'{}^\times)$ contains
that index-two subgroup by transitivity of norms.
It also contains every scalar cube.
Cubing is onto a group of order two, so these two
subgroups generate $k_n^\times$.
This proves (A.105.4), including the unit part.
On three-primary completions the norm is continuous
between compact local multiplicative completions.
Its image is closed and contains the dense image
of $k_n^\times$, so it is also onto there.
\end{proof}

The ray-modulus restriction in A.105 is an actual
hypothesis of its intersection proof: it uses
unramifiedness at three, and a ray modulus supported
there need not have that property.
The local conclusion (A.105.4) allows omission of
such a modulus for this particular norm equation.
It does not replace other independently prescribed
three-adic congruences.
At every permitted $q$ the $3^n$ old roots $\alpha_i$
are cubes, because they already have $3^{t-n}$th roots.
Thus their unequal cubic classes in A.103 disappear.
The separate target-unit equations (A.103.10) still
require a compatible choice of the repaired
representative; no recursive norm-chain extension
is asserted by this field-intersection result.

\noindent The exact certificate
\path{code/norm_deeper_root_certificate.py}, with output
\path{code/norm_deeper_root_certificate_output.txt}, checks
that the deeper translation acts through the stated
old index set, the literal arithmetic action identities,
the repeated valuation matrix, and the cyclic
coinvariant and index-two norm-quotient calculations.
